% !TEX root = ../manuscript.tex

\section{Offshore Texas Field Case and Geologic Uncertainty}

\subsection{Field Setting and Geologic Model}

We apply the workflow to an offshore Texas storage setting represented by a
three-dimensional geologic model with approximately two million grid cells.
The model contains a growth fault that offsets the storage interval and
partially offsets the top seal. This configuration permits migration both
along the fault and across juxtaposed units, making the spatially variable
fault properties central to predictions of containment.

\todo{Describe the storage interval, top-seal stratigraphy, model dimensions,
grid resolution, fault geometry, boundary conditions inherited from the field
model, and the evidence supporting the selected geologic interpretation. Add a
field-location and model-geometry figure.}

\subsection{Fault Discretization and Throw Windows}

The reservoir-scale fault representation contains 87 cells along strike and
six throw windows, denoted W1--W6. Each throw window has a different local
stratigraphic juxtaposition and therefore a distinct conditional distribution
of fault properties. The six window properties are assigned independently at
each along-strike location unless a sampled fault-wide or grouped case imposes
a shared low/high state. This $6\times87$ representation connects the
window-scale PREDICT libraries to the three-dimensional fault surface used by
the reservoir model.

\todo{Define the throw ranges associated with W1--W6 and show how the window
and along-strike indices map onto the reservoir grid.}

\subsection{Geologic Scenario Design}

Geologic uncertainty is the outer level of the hierarchy. We consider six
layer-thickness structures spanning low, medium, and high sand proportions,
each represented with uniform and nonuniform layer thicknesses. Each structure
is combined factorially with three faulting depths, three clay-volume
fractions in sand, and three clay-volume fractions in clay, producing

\begin{equation}
6\times3\times3\times3 = 162
\end{equation}

geologic scenarios. All scenarios receive equal treatment in the downstream
workflow; no clustering or representative-scenario reduction is applied at
this outer level.

\todo{Add the final parameter values, units, geologic rationale, and a table
that maps scenario identifiers to the complete factorial design.}
